\documentclass[a4paper]{scrartcl}
\usepackage[utf8]{inputenc}
\usepackage[english]{babel}
\usepackage{graphicx}
% SVG support: converted to vector PDF by Inkscape at compile time (needs -shell-escape, see latexmkrc).
% inkscapelatex=false keeps the diagram's own text and fonts instead of re-typesetting them with LaTeX.
\usepackage[inkscapelatex=false]{svg}
\usepackage{lastpage}
\usepackage{pgf}
\usepackage{wrapfig}
\usepackage{fancyvrb}
\usepackage{listings}
\usepackage{xcolor}
\usepackage{url}
\usepackage{lscape}
\usepackage[pass]{geometry}
\newcommand{\diagrampage}[3]{%
  \newgeometry{margin=10mm}%
  \begin{landscape}%
    \vspace*{\fill}
    \begin{figure}[h!]
      \centering
      \includesvg[width=\linewidth,height=0.78\paperwidth,keepaspectratio]{#1}
      \caption{#2}
      \label{#3}
    \end{figure}%
    \vspace*{\fill}
  \end{landscape}%
  \restoregeometry%
}

\lstdefinestyle{SQL}{
  language=SQL,
  morekeywords={RETURNS,TRIGGER,DECLARE,BEGIN,END,IF,THEN,RAISE,EXCEPTION,RETURN,LANGUAGE,FUNCTION,REPLACE,EXECUTE,BEFORE,EACH,ROW,INTEGER,LIMIT,INTO,INNER,JOIN},
  basicstyle=\ttfamily\footnotesize,
  keywordstyle=\color{blue!70!black}\bfseries,
  commentstyle=\color{green!45!black}\itshape,
  stringstyle=\color{red!65!black},
  frame=single,
  breaklines=true,
  columns=fullflexible,
  keepspaces=true,
  showstringspaces=false
}

% *** style added together with the filled-in sections ***
\lstdefinestyle{Java}{
  language=Java,
  basicstyle=\ttfamily\footnotesize,
  keywordstyle=\color{blue!70!black}\bfseries,
  commentstyle=\color{green!45!black}\itshape,
  stringstyle=\color{red!65!black},
  frame=single,
  breaklines=true,
  columns=fullflexible,
  keepspaces=true,
  showstringspaces=false
}

\lstdefinestyle{output}{
  basicstyle=\ttfamily\footnotesize,
  frame=single,
  breaklines=true,
  breakatwhitespace=false,
  columns=fullflexible,
  keepspaces=true,
  showstringspaces=false,
  literate={-}{{-}}1 {+}{{+}}1
}
\renewcommand{\_}{\textunderscore\allowbreak}
\emergencystretch=2em

\usepackage{fancyhdr}
\pagestyle{fancy}

% Create header and footer
\headheight 27pt
\pagestyle{fancyplain}
\lhead{\footnotesize{Data Storage Paradigms, IV1351}}
\chead{\footnotesize{Project task 3 report}}
\rhead{}
\lfoot{}
\cfoot{\thepage\ (\pageref{LastPage})}
\rfoot{}

\title{Data Storage Paradigms, IV1351}
\subtitle{Project task 3 report}
\author{}
\date{\today}

\begin{document}

\maketitle
\noindent\textbf{Project members:} \\ \hfill
[Mo Wang, mowang@ug.kth.se] \\ \hfill
% [Enter name, email] \\ \hfill
% [Enter name, email] \\ \hfill

\section*{Declaration:}

By submitting this assignment, it is hereby declared that all group members listed above have contributed to the solution. It is also declared that all project members fully understand all parts of the final solution and can explain it upon request.

It is furthermore declared that the solution below is a contribution by the project members only, and specifically that no part of the solution has been copied from any other source (except for lecture slides at the course IV1351), no part of the solution has been provided by someone not listed as a project member above, and no part of the solution has been generated by a system.

% \section*{Tips for Report Writing}
% \textbf{REMOVE THIS SECTION BEFORE SUBMITTING THE REPORT.}\\

% \noindent \textit{The target audience has exactly the same skills as the author, except they do not know anything at all about the specific application described in the report.} \\

% Consider the following:

% \begin{itemize}
%   \item \textbf{The report must be \textit{centered around the requirements}. Which are they (Introduction), how did you work to meet them (Method), what is the solution that meets them (Result), and how can you be sure they are met (Discussion). This is the IMRaD method.} The requirements on the Introduction, Method, Result and Discussion chapters are described below under each chapter.

%   \item Is spelling and grammar correct? Is spoken language avoided?

%   \item Does the report have a good structure with sections, subsections and paragraphs?

%   \item Is the text clarified with images and/or other figures, and with links to the code in your Git repository? Remember that all figures (images, tables, graphs, code listings, etc) shall be numbered and have a short explaining text.
% \end{itemize}

\section{Introduction}
\label{sec:introduction}

% \textbf{This chapter tells \textit{what} are you going to do.} 
% Explain the task and the requirements on the solution. It's important to clearly state the requirements. \textit{Also specify which other student you worked with when solving the tasks, or if you worked alone.}

This report presents the programmatic access of database which was designed in both project task 1 and 2. The database stores courses, course layout versions, course instances, planned teaching activities, employees and their teaching allocations at a university. The program is a part of course layout and teaching allocation's web site, which has certain functionalities of computing the teaching cost of a course isntance, modifying a course instance, allocate, deallocate teaching loads and adding teaching activities.

The database schema from project task 1 and 2 remains unchanged, since developing a database connection program doesn't need altering the underlying database. The database program is written in Java and connects to the database through JDBC. The program is divided into layers, following the layered MVC structure, where each responsibility, such as SQL scripts, business logic and user-facing interface, only stays in its dedicated layer. Additionally, the database program also handles transaction handling, by explicitly declaring beginning and end of a transaction, as well as rollbacks if failed and locking a table row using \texttt{FOR UPDATE} or \texttt{FOR NO KEY UPDATE <table\_name>}.

\subsection{Requirements of project task 3}
\label{sec:requirements}

The assignment is to develop part of course layout and teaching allocation's web site. a command line user interface is sufficient. The program is required to handle ACID transactions properly, which means autocommit must be turned off, instead the the program must call commit and rollback as required. Handling transactions properly also means that SELECT FOR UPDATE must be used when required.
Finally, you have to make sure the database contains sufficient data to check that all queries work as intended.

1. Compute the teaching cost (as planned and actually allocated) of a particular course
instance given in the current year. You may assume (or calculate) the average salary for a
teacher, for computation of the planned cost.

2. Modify the course instance: Increase the number of registered students by 100 to the
course instance you have selected at the previous step. Now compute the teaching cost for that
course instance again and see how the teaching cost is affected.

3. Allocate and deallocate teaching loads: Allocate and deallocate teaching activities for
various course instances for teachers. Remember that a teacher is not allowed to teach (or get involved in any teaching activities) in more than four course instances in a particular period, your program must check that this limit is not exceeded. You must demonstrate at least one allocation, one deallocation, and one case of throwing error for exceeding the limit.

4. Add a new teaching activity called "Exercise", this new activity must be associated with
at least one course layout/instance and must have one teacher allocated for the same. Also,
write a query to display the course layout/instance and allocation for the teacher which is
affected by this new activity.

\section{Literature Study}

% This chapter must prove that you collected sufficient knowledge before starting development, instead of just hacking away without knowing how to complete a task. State what you have read and briefly summarize what you have learned.

The knowledge required for project task 3 was collected from the following sources before writing the programmatic access for the database:

\begin{itemize}
  \item The project description~\cite{project} and the description of project task 3~\cite{task2} on the course pages in Canvas. They define the requirements summarized in Section~\ref{sec:requirements}.
  \item \textit{Tips and Tricks for Project Task 3}~\cite{tips2} and the \textit{Assessment Criteria For Seminar 3}~\cite{criteria2}. The assessment criteria were used as a checklist for the queries and for Chapter~\ref{sec:discussion}.
  \item Lecture on transaction
  \item Videos on JDBC introduction and Architecture and Design of a Database Application
  \item The course books~\cite{elmasri,silberschatz}, chapter 9 and 10, programmatic database access
  \item The PostgreSQL 16 documentation~\cite{postgresdocs}, chapter 13 on concurrency control and the reference page for \texttt{SELECT}, which specifies the locking clauses.
\end{itemize}

The project task 3 concerns OLTP instead of OLAP than in project task 2. The program is designed to handle simutanous transactional queries that touches a handful of rows. Additionally, the statements are relatively short without aggregate functions, and the difficulty relies more on performing concurrent actions than processing large volumes of data.

\subsection{Concurrent transaction}

A transaction is a sequence of statements that is considered as one atomic unit and satisfies the 4 ACID principles: atomicity, consistency, isolation and durability.

Atomicity is a property which a transaction must be fully completed or nothing at all, so a failure in the middle of a transaction must be rolled back. In SQL this is expressed by \texttt{COMMIT} and \texttt{ROLLBACK}.

Isolation means that each transaction produces the same behavior for sequential and concurrent execution. However, isolation is incomplete in practice since full isolation is expensive. In this database program, the only isolation anomaly that matters is lost update anomaly. This is because some operations fits into the read-modify-write pattern. This concerns having 2 transactions that both reads the row, compute a new value they read and write back to the row. If the first transaction reads but not yet written back to the row before the second transaction reading it, both transaction writes the same value, leading to a lost update. For preventing this, a row can be explicitly locked when reading, using \texttt{SELECT ... FOR UPDATE} or one of its weaker variants, which forces the second transaction to queue up until the first has committed. Additionally, the isolation level can also be reaised, which makes PostgreSQL detect the conflict and abort one of the transactions with a serialization failure that the program must then retry.

PostgreSQL offers several locking strengths. \texttt{FOR UPDATE} is the strongest and also conflicts with \texttt{FOR KEY SHARE}, which is the weak lock that PostgreSQL itself takes on a parent row when a row referencing it is inserted. \texttt{FOR NO KEY UPDATE} states that the row will be modified but its key will not, and therefore does not block such inserts, while still conflicting with other updates and with other row locks. Both prevent the lost update. The only difference is how much unrelated work they block.

\subsection{JDBC}

JDBC is the standard Java API for connection and interaction with relational databases. The API has several defined interfaces, including \texttt{Connection}, \texttt{Statement}, \texttt{PreparedStatement} and \texttt{ResultSet}. These are implemented by a driver specific to the database. This specific architecture decouples the application and the actual database-specific implementation. A connection is opened by \texttt{DriverManager.getConnection} with a URL, a user and a password.

There are 3 different API properties in JDBC. First, a JDBC connection is in auto-commit mode by default. This means that every statement is itself an implicit transaction. For multi-statement transactions, the autocommit must be set to false \texttt{setAutoCommit(false)} and explicit \texttt{commit()} and \texttt{rollback()} must be called. Second, \texttt{PreparedStatement}s can be compiled once with \texttt{?} placeholders and then executed with different parameter values. This offer both faster execution speed and native user input santization compared to raw SQL statement execution that takes time to compile and prone to SQL injection. Third, every JDBC method throws the checked \texttt{SQLException}, which is a vendor specific, database oriented exception. Letting it propagate would make the layers above the database aware of JDBC, so it has to be caught and wrapped.

\subsection{Architecture and Design of a Database Application}

The architecture of the database application consists of a layered MVC pattern, with View on the topmost layer, Controller, Model, Integration and Data layer. The Controller layer communicates with both Model and Integration, without any knowledge about View. Additionally, View has only knowledge about Controller. The integration layer, where database access objects (DAO) are located, are dumb and can only access the data layer and handle raw CRUD for the database. The principle is that each layer only calls the layer beneath it, with an exception to Controller layer that can also access Integration layer. This produces single direction of dependency and makes it possible to replace one layer without affecting other layers.

Each layer handles its intended and dedicated responsibility. For instance, Integration layer handles only the connection between actual data storage and model entities of the program. Also, Model layer handles the program entities, while Controller layer wraps the internal program buisness logic into system operations. Only the view layer contains logic about user interfaces.

The integration layer is implemented with the Data Access Object Pattern (DAO). A DAO holds the prepared statements and exposes methods that only handles raw data to database, without any buisness logic. The methods are limited to CRUD-operations, \texttt{create}, \texttt{read}, \texttt{update} and \texttt{delete}.

The data that cross layers are defined as Data Transfer Objects, which are immutable objects that only carries values. A DTO can be shaped as one query result, while a model object is an entity in buisness logic.

The videos on the database application architecture also made clear about the transaction boundaries. Since a transaction corresponds to one complete action, and the controller  is the layer that knows which operation is being performed, the controller is the only place that may call \texttt{commit} and \texttt{rollback}. A DAO sees only one statement of the operation and cannot know whether the operation as a whole has succeeded.

\section{Method}
% In the Method chapter of your report, mention which IDE(s) and other tool(s) you used and explain how you proceeded and reasoned when writing the program. Do not explain the result of each step you took, only explain the steps themselves.

In terms of method, the task is done first by writing the SQL queries for each operation. When the SQL queries produces expected behavior, they queries are translated as prepared statements. In the meantime, the layered architecture for the database access program is set up by putting placeholder classes for each layer: View, Controller, Model, Integration and Startup (Bootstrap). Once establishing the architecture skeleton, the code can be written from DAO integration layer up to the controller layer, while having a Basic printout view to verify the behavior of each Controller operation. A command line interpreter is written later.

\subsection{Tools}

For solving the task, the following tools were used:

The main database management system PostgreSQL 16 is used for developing analytical SQL scripts and improving their preformance. It is executed on Ubuntu 22 in Windows Subsystem for Linux (WSL), and all scripts were written as file and then executed using the psql command-line utility with the following command:

\texttt{sudo -u postgres psql -v ON\_ERROR\_STOP=1 -d course\_database -f "<file\_name>"}.

In this case, a query can be written and modified as a text file before feeding it into the database system.

For general development, Visual Studio Code was used to write the SQL script and this report, while Git and GitHub were used for version control for storing and backing up SQL queries. The report was written and typeset with \LaTeX.

For inserting additional data into the database, PostgreSQL function \texttt{generate\_series} is used. The number, often concatenated with additional prefixes and suffixes, for data values such as names, address and phone numbers.

In addition to the mentioned tools, the program was developed and executed with Maven, which handles Java multifile compilation and dependency manager, especially the PostgreSQL JDBC driver dependency. The program can simply be compiled, by cleaning the old artifact files if needed, with \texttt{mvn -o clean package} and executed using the command \texttt{mvn -o -q exec:java}. The source Java version is currently set to Java 17 in \texttt{pom.xml}, which is required because the DTOs are implemented as records.

The teacher's JDBC bank example was studied and run before any code was written, in order to see a working layered program against a database rather than only reading about one. It was kept in a separate folder of the repository and was not modified, so that the final program and the example can be compared.

\subsection{Writing each operation in SQL}

Every statement was first written in SQL in the folder \texttt{sql\_scripts} and manually run with \texttt{psql}, before translating it into a prepared statement in the program. For each operation, there is one script and each statement in the script approximately correspond to one DAO method. Since some parameters are unknown in the SQL scripts, most of them are given a default value, with comments that marks a specific parameter value to be user defined.

The reason for first developing each SQL statement before translating is that a wrong statement will become less obvious when executed in the program. In PostgreSQL CLI tool, miswritten SQL statements, such as a missing join condition or a forgotten multiplication factor is visible in the output table, since the whole table can be viewed. However, in the program, the same mistake only shows as a wrong number without too much information about surronding values. Troubleshooting a miswritten SQL statement would also be harder in the program, since the program produce less clues about the issue and the whole program has to be recompiled and re-executed. By developing SQL statements separately from the program, SQL statements that are translated into prepared statements can be guarenteed to be correct, and only the buisness logic needs to be verified when developing the program.

The translation from SQL script to prepared statements is mostly straightforward. Literal values that is a parameter value become \texttt{?} for prepared statement parameter. This is easy to distinguish, since parameters are wrapped with comments that indicates the literal as a specific parameter to be user defined. The transaction control statements \texttt{BEGIN}, \texttt{COMMIT}, \texttt{ROLLBACK} and \texttt{SAVEPOINT} are removed from the SQL script, because in JDBC those are expressed by \texttt{setAutoCommit(false)} and the corresponding methods on the connection rather than by SQL text. However, some parts that involves nested queries are removed and replaced since their values can be replaced with the results from earlier queries.

\subsection{Establishing layered architecture}

The layered architecture is established by creating one package/directory for each layer where each of them contains one or several placeholder class before writing and implementing any logic:

\begin{itemize}
  \item \texttt{se.kth.iv1351.startup}, which creates the objects and starts the program,
  \item \texttt{se.kth.iv1351.view}, the user interface,
  \item \texttt{se.kth.iv1351.controller}, which coordinates the operations,
  \item \texttt{se.kth.iv1351.model} and \texttt{se.kth.iv1351.dto}, the business logic and the objects transferred between layers,
  \item \texttt{se.kth.iv1351.integration}, the database access.
\end{itemize}

However, some packages require multiple classes, such as \texttt{integration} where \texttt{DatabaseConnection.java}, \texttt{AbstractDAO.java} and \texttt{TeachingActivityDAO.java} were created at once. The reason is that there are multiple database-facing classes in the Integration layer. By placing the packages first instead of letting them emerge later, a dependency that sits in a wrong direction or violate the layered MVC pattern becomes visible as a missing import. 

When the layered architecture scaffold is rolled out, there is a deliberately temporary construct for manually executing the program and testing. The View \texttt{BasicView} is a minimal View that calls different controller methods with hard-coded arguments and prints out the raw result.

\subsection{Deciding the model entities and DTO}

Both the model entities and DTOs are decided by modeling the functionality of the program. A class is only created when there is business logic tied to the entity, which produces the entities \texttt{CourseInstance} for increasing the number of students, \texttt{TeacherLoad} for the allocation rule and \texttt{TeachingCost} for turning hours and a salary into a cost. A candidate that turns out to be logicless, such as the employee read before the allocation check, is left as a DTO only.

% The DTOs were decided by writing the signature of each controller method before its body, since the return type of the method is the shape of the data the view needs. The DTOs therefore follow the results of the operations rather than the tables of the database, which is why \texttt{AllocationDTO} carries the teacher's name and the course code collected from several joined tables instead of the foreign keys that connect them.


\subsection{Writing each controller operation}

Every system operation in the program is written bottom-up, starting from the Integration layer by translating a SQL script into several SQL prepared statements and their DAO methods. There are multiple DAO implemented in the program, where each groups different related SQL statements.

Afterwards, the corresponding method for each system operation is written for the controller. The controller takes the arguments from the view and passes them down into different DAO methods and model entity objects. The body of each method is written as a sequence of DAO calls separated by the decisions taken between them. This made the transaction boundaries visible since the first DAO call starts a transaction implicitly and \texttt{commit} is the last statement before the return. In this case, the controller method chooses whether a lock should be used for a read, since controller methods contains the buisness logic and has the knowledge whether a lost-read anomaly could occur.

Each operation is verified against the database as soon as it compiled, by running the command and comparing the output with the result of the corresponding script in \texttt{sql\_scripts}. The changes made by the test runs are often undone, so that the figures reported in Section~\ref{sec:result} are reproducible from the data in the repository. Error handling is also tested in the same way, by giving the commands deliberately invalid arguments, which is how the two defects described in Chapter~\ref{sec:discussion} were found.


\section{Result}
\label{sec:result}
% In the Result chapter of your report, briefly explain the program and in particular explain ACID transaction handling. Include links to your git repository, and make sure the repository is public. Also include a printout of a sample run. The git repository must also contain the scripts that create the database and insert data. It shall be possible to test your solution by executing first the script that creates the database, then the script that inserts data, and finally execute your program.


\section{Discussion}
\label{sec:discussion}
% The Discussion chapter of your report must include a relevant and extensive evaluation of your program. Suggested assessment criteria are found in ownload seminar3-assessment-criteria.pdf, you do not have to cover them all. These same criteria will also be used to grade your project report.

% \section{Comments About the Course}

% \textbf{This section is optional, but please at least write approximately how much time you spent on the assignment}, including lectures, labs, tutorials and seminars. This is of great help for course evaluation.

% Also, any other comment(s) related to this course offering or to coming offerings is much appreciated. 

\section{References}
\label{sec:references}

% thebibliography is used instead of BibTeX, so that the report compiles with pdflatex
% alone and the reference list is kept in this single source file. Cite with \cite{key}.
\renewcommand{\refname}{}
\vspace{-2\baselineskip}
\begin{thebibliography}{9}

\bibitem{elmasri}
Ramez Elmasri and Shamkant B. Navathe.
\textit{Fundamentals of Database Systems}, 7th edition.
Pearson, 2016.

\bibitem{silberschatz}
Abraham Silberschatz, Henry F. Korth and S. Sudarshan.
\textit{Database System Concepts}, 7th edition (ISE).
McGraw-Hill, 2019.

\bibitem{project}
Leif Lindbäck and the IV1351 teaching staff.
\textit{Project: Course Layout and Teaching Load Allocations}.
Course page, Data Storage Paradigms IV1351, KTH, HT26.

% *** the three entries below were updated from the task 2 to the task 3 documents ***
\bibitem{task2}
\textit{Project Task 3, Programmatic Access to the Database}.
Course page, Data Storage Paradigms IV1351, KTH, HT26.

\bibitem{tips2}
\textit{Tips and Tricks for Project Task 3} (\texttt{seminar3-tips-and-tricks.pdf}),
including the JDBC bank example program.
Course material, Data Storage Paradigms IV1351, KTH, HT26.

\bibitem{criteria2}
\textit{Assessment Criteria For Seminar 3} (\texttt{seminar3-assessment-criteria.pdf}).
Course material, Data Storage Paradigms IV1351, KTH, HT26.

\bibitem{logicallecture}
Leif Lindbäck.
\textit{Logical and Physical Models}. Recorded lecture series,
Data Storage Paradigms IV1351, KTH, HT26.

\bibitem{sqllecture}
\textit{SQL, the Structured Query Language}. Lecture 3, Data Storage Paradigms IV1351, KTH, HT26.

\bibitem{normalisationlecture}
\textit{Normalisation}. Lecture 4, Data Storage Paradigms IV1351, KTH, HT26.

\bibitem{postgresdocs}
The PostgreSQL Global Development Group.
\textit{PostgreSQL 16 Documentation}.
\url{https://www.postgresql.org/docs/16/}.

\end{thebibliography}


\end{document}
